\documentclass[10pt,a4paper]{amsart}
\usepackage[margin=19mm]{geometry}
\usepackage{amsmath,amssymb,mathtools}
\usepackage[scr=rsfs,cal=euler]{mathalfa}
\usepackage{xfrac}
\usepackage[unicode,hidelinks,bookmarks=false]{hyperref}
\newcommand{\ie}{i.e.\ }
\newcommand{\cf}{cf.\ }
\newcommand{\resp}{respectively\ }
\newcommand{\Subsection}{\subsection}
\providecommand{\nobreakdash}{\nobreak\hbox{-}}
\setlength{\emergencystretch}{2em}
\multlinegap0pt
\usepackage{bm}
\usepackage{bbm}
\usepackage{dsfont}
\usepackage{xspace}
\usepackage{cancel}
\usepackage{mathtools}
\usepackage{amsthm}
\makeatletter
\@ifundefined{claim}{\newtheorem{claim}{Claim}}{}
\makeatother
\makeatletter
\expandafter\ifx\csname subalign\endcsname\relax
\newenvironment{subalign}{\subequations\align}{\endalign\endsubequations}
\fi
\makeatother
\title[The Fate of Crystalline Topological Phenomena in the Continu]{The Fate of Crystalline Topological Phenomena in the Continuum}
\author{Rajas Chari, Taylor L. Hughes}
\date{Source: arXiv:2608.11302v1 (CC BY 4.0)}
\begin{document}
\maketitle

\noindent\emph{Source and licence.} The Fate of Crystalline Topological Phenomena in the Continuum. Version arXiv:2608.11302v1 (CC BY 4.0), distributed under the Creative Commons Attribution 4.0 International licence, as recorded in the arXiv licence metadata for this exact version and shown on the abstract page https://arxiv.org/abs/2608.11302v1. Full rights notice: licenses/arxiv-2608.11302-CC-BY-4.0.txt.\par\smallskip

\noindent\textit{Editorial scope.} Counted unit: the Claim whose source label is \texttt{claim:degree\_dplus2\_acyclic} in the article (source0.tex lines 9272--9289, source label \texttt{claim:degree\_dplus2\_acyclic}), delivered together with the article's own complete original proof of it (source0.tex lines 9291--9414, one \texttt{proof} environment of 124 lines). No mathematics is written, repaired or re-derived: statement and proof are the article's own text line for line. Required context retained uncounted: none. Every definition, display and label of those lines is retained unchanged. Omitted from this package and not redistributed: 1--9271, 9290--9290, 9415--12913. Nothing inside a retained excerpt is omitted or altered: every retained body is delivered exactly as the article's lines are written, so no omission receipt of any kind accompanies this package. Citations: the retained excerpts carry no citation command at all, so no bibliography entry of the article is transcribed and no reference is redistributed. Reference adaptation: none was needed and none was made; every reference inside the delivered unit resolves inside it, so no unresolved reference or citation is produced. Printed slips kept literal: no printed word, symbol, hypothesis or number of the counted statement or of its proof was altered. Layout and class: the article's own class is \texttt{revtex4-2} with packages including booktabs, txfonts, fontenc, inputenc, hyperref, amsmath, color, graphicx, overpic, mathrsfs, bm, braket, acronym, comment; the wrapper is the lane's pinned \texttt{amsart} editorial wrapper at 10pt on A4, which restates the theorem-like environments the retained text needs and adds the article's own macro definitions that the retained text needs (none), with extra wrapper packages (bm, bbm, dsfont, xspace, cancel, mathtools). The delivered numbering is the unit's own. Assets: no retained span contains a figure environment, a graphics-inclusion command, a PDF-page inclusion, a table or a TikZ picture; no image file is copied into this package, the package declares no asset dependency and no image file is redistributed with it. Licence: version arXiv:2608.11302v1 of this article is distributed under the Creative Commons Attribution 4.0 International licence, as recorded in the arXiv licence metadata for this exact version and shown on the abstract page https://arxiv.org/abs/2608.11302v1. A full rights notice is delivered at licenses/arxiv-2608.11302-CC-BY-4.0.txt. Characters: every retained excerpt is pure ASCII, so the delivered engine, XeLaTeX with its default Latin Modern text font, reports no missing character.
\begin{claim}[Degree-$d+2$ acyclicity]
\label{claim:degree_dplus2_acyclic}
The degree-$d+2$ cohomology of $(\Lambda,D_{\lim})$ vanishes:
\begin{equation}
H^{d+2}(\Lambda,D_{\lim})=0.
\end{equation}
More precisely, if $d=2m+1$, then
\begin{equation}
H^{*}(\Lambda,D_{\lim})=0,
\end{equation}
whereas if $d=2m$, then
\begin{equation}
H^{*}(\Lambda,D_{\lim})
\cong
y_m\,\mathbb F_2[y_1^2,\dots,y_m^2].
\end{equation}
In particular, in the even-dimensional case the nonzero cohomology begins in degree $d$ and then occurs only in degrees $d+4k$, so it does not appear in degree $d+2$.
\end{claim}

\begin{proof}[\textnormal{\textbf{Proof of Claim~\ref{claim:degree_dplus2_acyclic}}}]
We use the decomposition
\begin{equation}
D_{\lim}=\delta+T,
\qquad
\delta=\sum_{t=1}^{m_o}z_t\partial_{y_t},
\qquad
T=z_0(1+\mathcal E),
\end{equation}
where $m_o=\lfloor (d-1)/2\rfloor$ and $m_e=\lfloor d/2\rfloor$. Recall that the essential difference between odd and even $d$ is that, when $d$ is even, the top even generator $y_{m_e}=e_d$ has no partner $z_{m_e}=e_{d+1}$.

\medskip

Define a decreasing filtration by powers of $z_0$:
\begin{equation}
F^p\Lambda:=z_0^p\Lambda,
\qquad
\Lambda=F^0\Lambda\supseteq F^1\Lambda\supseteq F^2\Lambda\supseteq \cdots .
\label{eq:z0_filtration}
\end{equation}
Since $\delta$ does not change the $z_0$-degree and $T$ increases it by one, we have
\begin{equation}
\delta(F^p\Lambda)\subseteq F^p\Lambda,
\qquad
T(F^p\Lambda)\subseteq F^{p+1}\Lambda\subseteq F^p\Lambda.
\end{equation}
Thus $D_{\lim}(F^p\Lambda)\subseteq F^p\Lambda$, so $(\Lambda,D_{\lim})$ is a filtered complex. In each fixed total degree this filtration is finite, and hence the associated spectral sequence converges to $H^*(\Lambda,D_{\lim})$.

\medskip

Let $A:=\Lambda/(z_0)$. Then $\Lambda\cong \bigoplus_{p\geq 0}z_0^p A$, and
\begin{equation}
\mathrm{gr}^p\Lambda:=\frac{F^p\Lambda}{F^{p+1}\Lambda}
\cong z_0^p A.
\label{eq:associated_graded_identification}
\end{equation}
\noindent The filtration above produces a spectral sequence whose first pages are the iterated cohomology groups
\begin{equation}
E_1\cong H_\delta(\Lambda),
\qquad
E_2\cong H_T\bigl(H_\delta(\Lambda)\bigr).
\label{eq:E1_E2_iterated}
\end{equation}
We now compute these pages, separating the odd and even cases.

\medskip

\noindent \textit{Odd-dimensional case.}
Suppose $d=2m+1$. Then $m_o=m_e=m$, and
$\Lambda\cong \mathbb F_2[z_0,y_1,z_1,\dots,y_m,z_m]$. In this case every $y_t$ is paired with a corresponding $z_t$. Since $\delta$ does not involve $z_0$, we first compute $H_\delta(A)$ and then adjoin $z_0$. For a single pair $(y,z)$, every polynomial can be written uniquely as
\begin{equation}
f(y,z)=f_0(y^2,z)+y f_1(y^2,z).
\end{equation}
Over $\mathbb F_2$, we have $\partial_y(y^{2k})=0$ and $\partial_y(y^{2k+1})=y^{2k}$, so
\begin{equation}
z\partial_y f=z f_1(y^2,z).
\end{equation}
Therefore $\ker(z\partial_y)=\mathbb F_2[y^2,z]$ and $\operatorname{Im}(z\partial_y)=z\cdot \mathbb F_2[y^2,z]$, which gives
\begin{equation}
H\bigl(\mathbb F_2[y,z],z\partial_y\bigr)
\cong
\mathbb F_2[y^2].
\end{equation}
Tensoring over the $m$ independent pairs $(y_t,z_t)$ gives
\begin{equation}
H_\delta(A)\cong \mathbb F_2[y_1^2,\dots,y_m^2],
\qquad
H_\delta(\Lambda)\cong \mathbb F_2[z_0,y_1^2,\dots,y_m^2].
\label{eq:Hdelta_result_odd}
\end{equation}

Now compute the induced $T$-cohomology. A class in $H_\delta(\Lambda)$ is represented by $z_0^p f(y_1^2,\dots,y_m^2)$. Since each $y_t$ appears only with even exponent in $f$, the Euler operator satisfies $\mathcal E(f)\equiv 0\pmod 2$. Hence
\begin{equation}
T\bigl(z_0^p f\bigr)
=
z_0(1+\mathcal E)\bigl(z_0^p f\bigr)
=
(1+p)z_0^{p+1}f.
\end{equation}
Thus the induced differential sends even powers of $z_0$ isomorphically onto the odd powers and vanishes on odd powers. Equivalently, if
$E_1^{\mathrm{even}}=\mathbb F_2[z_0^2,y_1^2,\dots,y_m^2]$ and
$E_1^{\mathrm{odd}}=z_0\mathbb F_2[z_0^2,y_1^2,\dots,y_m^2]$, then the induced differential is multiplication by $z_0$ from $E_1^{\mathrm{even}}$ to $E_1^{\mathrm{odd}}$ and vanishes on $E_1^{\mathrm{odd}}$. Since $E_1^{\mathrm{even}}$ has no zero divisors, this gives
\begin{equation}
E_2=H_T\bigl(H_\delta(\Lambda)\bigr)=0.
\end{equation}
Therefore the spectral sequence collapses to zero, and
\begin{equation}
H^*(\Lambda,D_{\lim})=0
\end{equation}
when $d$ is odd. In particular, $H^{d+2}(\Lambda,D_{\lim})=0$.

\medskip

\noindent \textit{Even-dimensional case.}
Now suppose $d=2m$. Then $m_e=m$ and $m_o=m-1$, so
\begin{equation}
\Lambda\cong \mathbb F_2[z_0,y_1,z_1,\dots,z_{m-1},y_m].
\end{equation}
The odd-dimensional computation fails in exactly one place: the top even generator $y_m=e_d$ has no partner $z_m=e_{d+1}$. Thus $\delta$ pairs $y_1,\dots,y_{m-1}$ with $z_1,\dots,z_{m-1}$, but leaves $y_m$ unpaired.

\noindent Repeating the same $\delta$-cohomology computation gives
\begin{equation}
H_\delta(\Lambda)
\cong
\mathbb F_2[z_0,y_1^2,\dots,y_{m-1}^2,y_m].
\label{eq:Hdelta_result_even}
\end{equation}
The induced $T$-differential is still $z_0(1+\mathcal E)$ on this ring. The same parity computation as above now leaves precisely the odd powers of the unpaired generator $y_m$. Thus
\begin{equation}
H^*(\Lambda,D_{\lim})
\cong
y_m\,\mathbb F_2[y_1^2,\dots,y_{m-1}^2,y_m^2].
\label{eq:even_cohomology_result}
\end{equation}
This is the precise failure of full acyclicity in even dimension.

\bigskip 

\noindent However, this failure does not affect the degree needed here. Since $|y_m|=d$ and the remaining generators in \eqref{eq:even_cohomology_result} have degrees divisible by $4$, the nonzero cohomology occurs only in degrees $d,d+4,d+8,\dots$. In particular, there is no cohomology in degree $d+2$. Therefore
\begin{equation}
H^{d+2}(\Lambda,D_{\lim})=0
\end{equation}
also in the even-dimensional case. Combining the odd and even cases proves the claim.
\end{proof}

\end{document}
